%% ============================================================
%%  Estado del arte
%%  Duenio inicial: claude-1
%%  Reglas: ver ../../CLAUDE.md
%%  No editar si TASKS.md marca este archivo IN_PROGRESS por otro agente.
%% ============================================================

\chapter{Estado del arte}
\label{cap:estado_arte}

\section{Obtención del corpus}

La revisión se condujo bajo el protocolo PRISMA 2020, elegido porque documenta cada decisión de selección y permite que un tercero reconstruya el corpus a partir del mismo registro. De la pregunta de investigación presentada en el capítulo anterior se derivó una cadena de búsqueda con cuatro bloques de términos, uno por cada componente de la pregunta, unidos entre sí con el operador \texttt{AND} y con sinónimos enlazados por \texttt{OR} dentro de cada bloque. Esa cadena se ejecutó sobre Scopus, Web of Science y PubMed, con los resultados que recoge la tabla \ref{tab:bases}, y se complementó con búsquedas por citas y registros.

Los criterios de inclusión restringieron la búsqueda a publicaciones entre 2022 y 2026, en inglés, y a artículos de investigación primaria o trabajos de conferencia, sin restricción de tipo documental en PubMed para no perder literatura técnica no clínica. El cribado posterior se llevó en Zotero, que permitió depurar duplicados y descartar registros sin identificador auditable o sin texto completo disponible. De ahí salieron las 25 referencias que componen el corpus de este informe, todas leídas de forma completa y registradas en una matriz de extracción con su descripción, muestra, metodología, resultados y conclusión, junto con la sección de la tesis a la que aporta cada una.

\begin{table}[htbp]
\centering
\caption{Registros identificados por base de datos.}
\label{tab:bases}
\begin{tabular}{lr}
\toprule
\textbf{Base de datos} & \textbf{Registros} \\
\midrule
Scopus & 191 \\
Web of Science & 100 \\
PubMed & 28 \\
\midrule
\textbf{Total identificado} & \textbf{319} \\
\bottomrule
\end{tabular}
\end{table}

\section{Trabajos similares}

El contraste se organiza por familias y no trabajo por trabajo, porque investigaciones que a primera vista no se parecen terminan compartiendo la misma limitación estructural. Lo que define cada familia es el tipo de objeto que entrega como explicación y el momento del ciclo de vida del modelo en que ese objeto se produce.

\subsection{Explicabilidad posterior al entrenamiento y su fiabilidad}

La primera familia calcula relevancia sobre un modelo ya entrenado, que es la práctica dominante cuando hay que explicar un sistema opaco sin tocarlo. \textcite{chefer2021transformer} propusieron calcularla mediante descomposición de Taylor propagada por capas, incluidas las conexiones residuales, y su método supera tanto a los mapas de atención crudos como a la propagación heurística. La dificultad aparece al comparar métodos entre sí. Sobre un clasificador Transformer de texto médico, SHAP y gradientes integrados alcanzan una fidelidad prácticamente igual, en torno a 0,22, y sin embargo solo correlacionan de forma moderada \parencite{medtext2026attribution}, es decir, dos explicaciones igualmente fieles ordenan las variables de manera distinta. El problema se generaliza en la evaluación sistemática de \textcite{vitexplain2024evaluation}, donde métricas que dicen medir la misma propiedad convergen mínimamente entre sí, con la consecuencia incómoda de que el orden de mérito entre métodos depende de cuál métrica se haya elegido para construirlo.

A esa inconsistencia se suma un cuerpo de trabajos que mide directamente cuánto se puede confiar en una explicación, y cuyas conclusiones condicionan el diseño de la propuesta más que cualquier otra familia. Ya se mencionó el ataque de \textcite{slack2020fooling} sobre los métodos basados en perturbación y el caso de aprendizaje por atajo que documentaron \textcite{lapuschkin2019unmasking}. Un tercer resultado, más reciente, agrega una tensión que el diseño tiene que administrar: al reentrenar redes convolucionales y Transformers de visión con siete semillas sobre la misma tarea, la arquitectura de atención produce explicaciones más parecidas entre corridas, pero borrar las regiones que destaca apenas altera su confianza \parencite{keratoconus2026instability}. Estabilidad y fidelidad, entonces, no son la misma propiedad y pueden moverse en direcciones opuestas. El límite común de toda esta familia es el tipo de objeto que produce, una importancia por variable o un mapa continuo, sobre el cual no se puede formular una hipótesis enumerable.

\subsection{Estructura relacional aprendida, impuesta o intervenida}

La segunda familia sí produce estructura, aunque la obtiene por vías que la propuesta no puede seguir. Una parte la aprende como componente del modelo: \textcite{cai2024msgnet} combinan descomposición en frecuencia con convolución de grafo adaptativa y exponen un grafo interpretable por escala temporal. Otra parte la recibe de antemano, como en el Transformer de grafo de \textcite{zhao2026causalguided}, que construye un grafo de propagación de fallos mediante cribado de causalidad de Granger y lo usa para condicionar el cálculo de la atención, de modo que el análisis de interpretabilidad termina realizándose sobre una topología que el propio modelo no puede poner a prueba. El trabajo de \textcite{pipeline2026causalgraph} es el más sólido del grupo en cuanto a validación, porque contrasta el grafo de dependencias condicionales que aprende contra la topología física real de la instalación que modela, y también el que mejor expone la dificultad del problema aquí planteado: esa validación es posible porque existe una referencia externa, y entre las variables del caso de estudio no la hay. Cierra el grupo la propuesta de \textcite{rulexai2024events}, que destila un clasificador en un conjunto compacto de reglas y lo evalúa a la vez por fidelidad y por tamaño, criterio doble que este trabajo adopta.

Un caso aparte, y el más cercano a la propuesta, es el de \textcite{mechinterp2025ts}, que traslada técnicas de interpretabilidad mecanicista desde el procesamiento de lenguaje a un Transformer de clasificación de series temporales y construye grafos causales del flujo interno de información entre cabezas de atención. El objeto que obtiene es exactamente el que aquí se persigue. La diferencia está en el precio: para conseguirlo interviene el modelo mediante parcheo de activaciones y necesita acceso a sus estados internos, lo que descarta su aplicación sobre un modelo ya desplegado del que solo se dispone de la matriz de atención. Ese contraste define por oposición el alcance de la propuesta, que renuncia a intervenir el modelo y asume el costo de tener que validar de otra manera.

\subsection{Arquitecturas afines y aplicaciones de dominio}

La tercera familia no aborda interpretabilidad, sino dónde conviene aplicar la atención y qué se gana con ello. Junto al resultado ya citado de \textcite{liu2024itransformer}, la arquitectura de \textcite{sentinel2025multipatch} combina atención de canal con atención temporal mediante parches múltiples y alcanza resultados comparables o mejores que el estado del arte. En sentido contrario apunta el trabajo de \textcite{patchmlp2025unlocking}, para quienes buena parte del mérito atribuido a los Transformers en pronóstico de largo plazo procede del mecanismo de parcheo y no de la atención entre variables, afirmación que sostienen mostrando que un perceptrón multicapa sobre parches iguala a arquitecturas bastante más complejas. La familia respalda la condición que hace posible este trabajo, la existencia de una matriz de atención entre variables, sin resolver de dónde viene la ganancia predictiva.

Las aplicaciones a dominios cercanos completan el panorama y muestran el hueco con más claridad que cualquier argumento teórico. El modelo de \textcite{convlstmgcn2026vegetation} predice mapas de índices de vegetación con error bajo y transfiere entre sensores, sin auditar en ningún momento la estructura que aprendió. \textcite{teleconnections2026causal} insertan un grafo acíclico dirigido con sentido físico en el espacio latente de un autocodificador variacional y validan el resultado contra experimentos numéricos, aunque la estructura causal en su caso viene impuesta y no extraída. \textcite{spatiotemporal2024e2e} desarrollan explicaciones agnósticas al modelo para redes de sensores espacio-temporales y, de paso, muestran que las herramientas pensadas para datos tabulares o de imagen resultan inefectivas en ese contexto. El caso más llamativo es el de \textcite{rotorcraft2026vortexring}, donde la atención de un Transformer de visión converge por sí sola sobre una banda de frecuencia que la teoría aerodinámica predice de forma independiente. Es el único trabajo del corpus donde la atención recupera un mecanismo verificable por fuera del modelo, y precisamente por eso resulta útil: ocurrió una vez, por coincidencia con una magnitud ya conocida, no por un procedimiento que pueda repetirse en otro problema.

\section{Contribución de la propuesta}

La tabla \ref{tab:familias} resume el contraste. Leídas en conjunto, las tres familias dibujan un panorama asimétrico, donde obtener un objeto relacional está resuelto por varias vías y decidir si ese objeto merece crédito no lo está.

Lo que falta es un procedimiento para falsar una estructura leída, es decir, ni aprendida como parámetro, ni impuesta como conocimiento previo, ni obtenida interviniendo el modelo. Los trabajos que evalúan explicaciones lo hacen sobre atribuciones por variable, con métricas que no convergen entre sí \parencite{vitexplain2024evaluation,medtext2026attribution}, y no sobre las relaciones de un grafo. Los que aplican estas arquitecturas a dominios afines predicen bien sin auditar lo que aprendieron \parencite{convlstmgcn2026vegetation,spatiotemporal2024e2e}. Y el único caso donde la atención resulta informativa por sí sola llegó ahí por coincidencia con un mecanismo conocido de antemano \parencite{rotorcraft2026vortexring}.

La propuesta ocupa ese espacio combinando cinco elementos que el corpus revisado presenta por separado y nunca juntos: extracción posterior al entrenamiento sin reentrenar ni intervenir el modelo, consenso frente a una hipótesis nula explícita, medición de fidelidad relación por relación, contraste entre paradigmas que no comparten supuestos, y una taxonomía que separa relación genuina de atajo predictivo, artefacto y ruido. Cada uno tiene precedente aislado en la literatura; lo que no tiene precedente es la composición.

{\small
\begin{longtable}{P{0.22}P{0.22}P{0.28}P{0.28}}
\caption{Familias del estado del arte, con su aporte y su límite.}\label{tab:familias}\\
\toprule
\textbf{Familia} & \textbf{Trabajos} & \textbf{A favor} & \textbf{En contra}\\
\midrule
\endfirsthead
\toprule
\textbf{Familia} & \textbf{Trabajos} & \textbf{A favor} & \textbf{En contra}\\
\midrule
\endhead
Explicabilidad posterior al entrenamiento y su fiabilidad & \parencite{chefer2021transformer,medtext2026attribution,vitexplain2024evaluation,slack2020fooling,lapuschkin2019unmasking,keratoconus2026instability} & Se aplican sin reentrenar y acotan con evidencia experimental cuánto se puede confiar en una explicación & Entregan importancia por variable o mapa continuo; las métricas que las comparan no convergen entre sí\\
Estructura aprendida, impuesta o intervenida & \parencite{cai2024msgnet,zhao2026causalguided,pipeline2026causalgraph,rulexai2024events,mechinterp2025ts} & Producen estructura explícita; un caso la contrasta contra topología física y otro obtiene un grafo causal verificable & Auditarla exige reentrenar, o es supuesto previo sin hipótesis nula, o requiere intervenir las activaciones\\
Arquitecturas afines y aplicaciones de dominio & \parencite{liu2024itransformer,sentinel2025multipatch,patchmlp2025unlocking,convlstmgcn2026vegetation,teleconnections2026causal,spatiotemporal2024e2e,rotorcraft2026vortexring} & Sostienen que atender entre variables es efectivo y muestran viabilidad en dominios cercanos & No zanjan si la ganancia viene de la atención o del parcheo, y no auditan la estructura aprendida\\
\textbf{Propuesta} & \multicolumn{3}{p{0.80\linewidth}}{Extracción posterior al entrenamiento sin intervenir el modelo, con criterio de validación acumulativo, hipótesis nula explícita, contraste entre paradigmas y taxonomía diagnóstica}\\
\bottomrule
\end{longtable}
}
